\section{Niveau 3}

\begin{UPSTIinfor}{Imbriquer ou enchaîner}
    Deux structures conditionnelles sont dites \textbf{enchaînées} lorsqu'elles se succèdent au même niveau, avec \texttt{else if} :
    \begin{lstlisting}[language=C]
if (cond1) { ... }
else if (cond2) { ... }
else { ... }
    \end{lstlisting}
    Elles sont dites \textbf{imbriquées} (ou \og nichées \fg) lorsqu'un \texttt{if} complet se trouve \textbf{à l'intérieur} du bloc d'un autre \texttt{if} :
    \begin{lstlisting}[language=C]
if (cond1)
{
    if (cond2) { ... }
    else { ... }
}
else
{
    ...
}
    \end{lstlisting}
    Dans ce second cas, \texttt{cond2} n'est testée que si \texttt{cond1} est déjà vraie : les deux conditions ne sont pas au même \og niveau \fg.
\end{UPSTIinfor}

\subsection{Imbrication, ordre et bornes}

\begin{UPSTIexercice}{Imbriqué ou enchaîné : deux écritures pour classer un nombre}
    On veut classer un entier \texttt{x} en trois catégories : \og Négatif \fg{} (\texttt{x < 0}), \og Petit positif \fg{} (\texttt{0 <= x < 10}), \og Grand positif \fg{} (\texttt{x >= 10}). Deux versions sont proposées.

    \noindent
    \begin{minipage}[t]{0.48\linewidth}
        \textbf{Version A (imbriquée)}
        \begin{lstlisting}[language=C]
if (x >= 0)
{
    if (x < 10)
    {
        printf("Petit positif\n");
    }
    else
    {
        printf("Grand positif\n");
    }
}
else
{
    printf("Negatif\n");
}
        \end{lstlisting}
    \end{minipage}\hfill
    \begin{minipage}[t]{0.48\linewidth}
        \textbf{Version B (enchaînée)}
        \begin{lstlisting}[language=C]
if (x >= 0 && x < 10)
{
    printf("Petit positif\n");
}
else if (x >= 10)
{
    printf("Grand positif\n");
}
else
{
    printf("Negatif\n");
}
        \end{lstlisting}
    \end{minipage}

    \UPSTIquestion{Dans la version A, \textbf{entourer} le \texttt{if} imbriqué (celui qui se trouve à l'intérieur d'un autre bloc).}
    \UPSTIquestion{Tracer l'exécution des deux versions pour \texttt{x = -3}, \texttt{x = 4} et \texttt{x = 15}. Affichent-elles le même résultat dans les trois cas ?}
    \UPSTIquestion{Dans la version A, le test \texttt{x < 10} n'apparaît que dans le cas où \texttt{x >= 0} est déjà vrai. Pourquoi la version B doit-elle, elle, écrire explicitement \texttt{x >= 0 \&\& x < 10} et ne peut-elle pas se contenter de \texttt{x < 10} ?}
\end{UPSTIexercice}

\begin{UPSTIprofOnlyEnv}
    \begin{UPSTIcorrectionP}{Imbriqué ou enchaîné : deux écritures pour classer un nombre}
        Le \texttt{if (x < 10) ... else ...} imbriqué dans la version A est celui qui se trouve à l'intérieur du bloc du premier \texttt{if}.

        Trace :
        \begin{itemize}
            \item \texttt{x = -3} : A $\to$ \texttt{x >= 0} faux $\to$ \og Negatif \fg. B $\to$ première condition fausse, \texttt{x >= 10} fausse $\to$ \og Negatif \fg. \textbf{Même résultat.}
            \item \texttt{x = 4} : A $\to$ \texttt{x >= 0} vrai, \texttt{x < 10} vrai $\to$ \og Petit positif \fg. B $\to$ \texttt{x >= 0 \&\& x < 10} vrai $\to$ \og Petit positif \fg. \textbf{Même résultat.}
            \item \texttt{x = 15} : A $\to$ \texttt{x >= 0} vrai, \texttt{x < 10} faux $\to$ \og Grand positif \fg. B $\to$ première condition fausse (15 n'est pas $<$10), \texttt{x >= 10} vrai $\to$ \og Grand positif \fg. \textbf{Même résultat.}
        \end{itemize}
        Les deux versions sont donc équivalentes.

        Dans la version A, le test interne \texttt{x < 10} hérite implicitement du fait que l'on est déjà dans le cas \texttt{x >= 0} : c'est la structure du code (l'imbrication) qui porte cette information. Dans la version B, les deux tests sont au même niveau, indépendants l'un de l'autre : rien ne garantit qu'on est passé par une vérification de \texttt{x >= 0} avant de tester \texttt{x < 10}. Il faut donc \textbf{répéter explicitement} la condition \texttt{x >= 0} dans le \texttt{\&\&} pour obtenir le même comportement.
    \end{UPSTIcorrectionP}
\end{UPSTIprofOnlyEnv}


\begin{UPSTIexercice}{À la limite}
    Une billetterie applique la règle suivante : \og le tarif enfant s'applique aux personnes dont l'âge est inférieur ou égal à 12 ans \fg. Un développeur a écrit :
    \begin{lstlisting}[language=C]
if (age < 12)
{
    printf("Tarif enfant\n");
}
else
{
    printf("Tarif normal\n");
}
    \end{lstlisting}
    \UPSTIquestion{Tracer l'exécution de ce code pour \texttt{age = 11}, \texttt{age = 12} et \texttt{age = 13}.}
    \UPSTIquestion{Ce code respecte-t-il exactement la règle énoncée ? Justifier en identifiant précisément la valeur d'âge pour laquelle le résultat n'est pas celui attendu.}
    \UPSTIquestion{Corriger le code.}
\end{UPSTIexercice}

\begin{UPSTIprofOnlyEnv}
    \begin{UPSTIcorrectionP}{À la limite}
        \texttt{age = 11} : \texttt{11 < 12} vrai $\to$ \og Tarif enfant \fg{} (conforme).
        \texttt{age = 12} : \texttt{12 < 12} \textbf{faux} $\to$ \og Tarif normal \fg.
        \texttt{age = 13} : \texttt{13 < 12} faux $\to$ \og Tarif normal \fg{} (conforme).

        Le code ne respecte \textbf{pas} exactement la règle : pour \texttt{age = 12}, la règle demande \og Tarif enfant \fg{} (12 ans \textbf{inclus}), mais le code, avec un \texttt{<} strict, affiche \og Tarif normal \fg. C'est une erreur classique \og d'un cran \fg{} (\emph{off-by-one}) à la frontière.

        Correction :
        \begin{lstlisting}[language=C]
if (age <= 12)
{
    printf("Tarif enfant\n");
}
else
{
    printf("Tarif normal\n");
}
        \end{lstlisting}
    \end{UPSTIcorrectionP}
\end{UPSTIprofOnlyEnv}


\begin{UPSTIexercice}{Une branche qu'on n'atteindra jamais}
    Un radar routier affiche un message selon la vitesse mesurée \texttt{vitesse} (en km/h) :
    \begin{lstlisting}[language=C]
if (vitesse >= 0)
{
    printf("Vitesse normale\n");
}
else if (vitesse > 130)
{
    printf("Exces de vitesse\n");
}
else
{
    printf("Vitesse negative (capteur en panne)\n");
}
    \end{lstlisting}
    \UPSTIquestion{Tracer l'exécution pour \texttt{vitesse = 150}. Le message \og Exces de vitesse \fg{} s'affiche-t-il ?}
    \UPSTIquestion{Une des trois branches de ce code ne s'exécutera \textbf{jamais}, quelle que soit la valeur de \texttt{vitesse}. Laquelle, et pourquoi ?}
    \UPSTIquestion{Corriger l'ordre des conditions pour que le programme se comporte comme voulu (excès de vitesse au-delà de 130, vitesse normale entre 0 et 130 inclus, capteur en panne en dessous de 0).}
\end{UPSTIexercice}

\begin{UPSTIprofOnlyEnv}
    \begin{UPSTIcorrectionP}{Une branche qu'on n'atteindra jamais}
        Pour \texttt{vitesse = 150} : \texttt{vitesse >= 0} est vraie $\to$ on affiche directement \og Vitesse normale \fg. Le \texttt{else if (vitesse > 130)} n'est même pas testé.

        La branche \og Exces de vitesse \fg{} est \textbf{inatteignable} : toute valeur qui vérifie \texttt{vitesse > 130} vérifie \emph{aussi} \texttt{vitesse >= 0} (puisque 130 $>$ 0). La première condition, trop large, \og capture \fg{} systématiquement les cas visés par la seconde, qui n'est donc jamais testée.

        Correction : il faut tester la condition la plus restrictive en premier.
        \begin{lstlisting}[language=C]
if (vitesse > 130)
{
    printf("Exces de vitesse\n");
}
else if (vitesse >= 0)
{
    printf("Vitesse normale\n");
}
else
{
    printf("Vitesse negative (capteur en panne)\n");
}
        \end{lstlisting}
    \end{UPSTIcorrectionP}
\end{UPSTIprofOnlyEnv}


\subsection{Équivalences et pièges logiques}

\begin{UPSTIexercice}{Deux écritures, un seul résultat ?}
    On considère deux entiers \texttt{a} et \texttt{b}, et les trois expressions booléennes suivantes :
    \begin{itemize}
        \item E1 : \lstinline[language=C]{!(a > 0 && b > 0)}
        \item E2 : \lstinline[language=C]{!(a > 0) || !(b > 0)}
        \item E3 : \lstinline[language=C]{!(a > 0) && !(b > 0)}
    \end{itemize}
    \UPSTIquestion{Compléter le tableau ci-dessous pour les quatre couples de valeurs indiqués.}

    \begin{center}
        \begin{tabular}{|c|c|c|c|c|c|}
            \hline
            \texttt{a} & \texttt{b} & \texttt{a > 0} & \texttt{b > 0} & E1, E2, E3 (vrai/faux) \\
            \hline
            5          & 3          &                &                &                        \\[0.5cm]
            \hline
            5          & -3         &                &                &                        \\[0.5cm]
            \hline
            -5         & 3          &                &                &                        \\[0.5cm]
            \hline
            -5         & -3         &                &                &                        \\[0.5cm]
            \hline
        \end{tabular}
    \end{center}
    \UPSTIquestion{Laquelle de E2 ou E3 est réellement \textbf{équivalente} à E1 pour tous les couples \texttt{(a, b)}, pas seulement ceux du tableau ? Justifier sans recalculer, en raisonnant sur le sens de chaque expression.}
    \UPSTIquestion{À quelle expression plus simple l'expression restante (celle qui n'est \textbf{pas} équivalente à E1) correspond-elle ?}
\end{UPSTIexercice}

\begin{UPSTIprofOnlyEnv}
    \begin{UPSTIcorrectionP}{Deux écritures, un seul résultat ?}
        \begin{center}
            \begin{tabular}{|c|c|c|c|c|c|c|c|}
                \hline
                \texttt{a} & \texttt{b} & \texttt{a>0} & \texttt{b>0} & E1 & E2 & E3 \\
                \hline
                5          & 3          & V            & V            & F  & F  & F  \\
                \hline
                5          & -3         & V            & F            & V  & V  & F  \\
                \hline
                -5         & 3          & F            & V            & V  & V  & F  \\
                \hline
                -5         & -3         & F            & F            & V  & V  & V  \\
                \hline
            \end{tabular}
        \end{center}

        E1 et E2 sont \textbf{équivalentes} sur les quatre lignes, et le sont en réalité pour \textbf{tout} couple \texttt{(a, b)} : \og il est faux que les deux soient positifs \fg{} (E1) revient exactement à dire \og l'un au moins des deux n'est pas positif \fg{} (E2). Nier un \og ET \fg{} revient à \og OU \fg{} des négations : dès qu'une des deux conditions d'origine est fausse, la négation du \texttt{ET} devient vraie, et c'est justement ce que teste le \texttt{OU} des négations.

        E3 correspond en réalité à \lstinline[language=C]{!(a > 0 || b > 0)} (\og aucun des deux n'est positif \fg), ce qui est une affirmation plus forte que E1 (\og au moins un des deux n'est pas positif \fg) : E3 est vraie uniquement quand les deux sont non positifs à la fois, comme le montre la dernière ligne du tableau.
    \end{UPSTIcorrectionP}
\end{UPSTIprofOnlyEnv}


\begin{UPSTIexercice}{Ça compile, mais est-ce juste ?}
    On exécute le programme suivant, qui a compilé sans aucune erreur ni avertissement.
    \begin{lstlisting}[language=C]
int x = 3;

if (x = 5)
{
    printf("x vaut 5\n");
}
else
{
    printf("x ne vaut pas 5\n");
}

printf("Valeur finale de x : %d\n", x);
    \end{lstlisting}
    \UPSTIquestion{Quel est l'opérateur utilisé dans \texttt{if (x = 5)} ? Est-ce une comparaison ou une affectation ?}
    \UPSTIquestion{Donner la valeur de \texttt{x} juste après l'exécution de la ligne \texttt{if (x = 5)}, puis la valeur utilisée par le \texttt{if} pour choisir sa branche.}
    \UPSTIquestion{Prédire, ligne par ligne, tout ce que ce programme affiche.}
    \UPSTIquestion{Ce programme se comporte-t-il comme si son auteur avait voulu \emph{tester} si \texttt{x} valait 5 ? Corriger l'erreur.}
\end{UPSTIexercice}

\begin{UPSTIprofOnlyEnv}
    \begin{UPSTIcorrectionP}{Ça compile, mais est-ce juste ?}
        \texttt{x = 5} est une \textbf{affectation} (un seul signe \texttt{=}), pas une comparaison : elle range 5 dans \texttt{x}, puis toute l'expression \texttt{(x = 5)} \og vaut \fg{} la valeur affectée, ici 5.

        Après cette ligne, \texttt{x} vaut \textbf{5}. Le \texttt{if} regarde ensuite la valeur de l'expression, c'est-à-dire \textbf{5}. En C, toute valeur non nulle est considérée comme vraie : le \texttt{if} prend donc systématiquement la branche du \texttt{if}, \textbf{quelle qu'ait été la valeur initiale de \texttt{x}}.

        Affichage :
        \begin{lstlisting}[language=bash,style=console]
x vaut 5
Valeur finale de x : 5
        \end{lstlisting}

        Non : le code \textbf{compile et s'exécute}, mais il ne teste jamais réellement si \texttt{x} valait 5 avant cette ligne --- il l'y force. C'est un exemple de code syntaxiquement correct mais logiquement faux (voir aussi la mise en garde du cours sur \texttt{=} et \texttt{==}). Correction :
        \begin{lstlisting}[language=C]
if (x == 5)
        \end{lstlisting}
    \end{UPSTIcorrectionP}
\end{UPSTIprofOnlyEnv}


\subsection{Synthèse}

\begin{UPSTIexercice}{Tarification d'un parking (exercice de synthèse)}
    On souhaite programmer l'affichage du tarif d'un parking, à partir de deux informations : le type de véhicule (\texttt{char type}, valant \texttt{'M'} pour moto ou \texttt{'V'} pour voiture) et la durée de stationnement en minutes (\texttt{int dureeMinutes}).

    Règles à respecter :
    \begin{itemize}
        \item Pour une \textbf{moto} : gratuit si \texttt{dureeMinutes <= 15} ; tarif de \texttt{2} sinon.
        \item Pour une \textbf{voiture} : gratuit si \texttt{dureeMinutes <= 15} ; tarif de \texttt{3} si \texttt{dureeMinutes} est compris entre 16 et 60 minutes (bornes incluses) ; tarif de \texttt{6} au-delà de 60 minutes.
        \item Si \texttt{type} ne vaut ni \texttt{'M'} ni \texttt{'V'}, afficher \og Type de vehicule inconnu \fg.
    \end{itemize}
    On affichera uniquement \texttt{"Tarif : "} suivi du montant (par exemple \texttt{"Tarif : 3"}), ou le message d'erreur.

    \UPSTIquestion{Écrire le pseudo-code complet de ce programme. Vous devrez nécessairement \textbf{imbriquer} une structure conditionnelle à l'intérieur d'une autre.}

    \UPSTIquestion{Traduire ce pseudo-code en langage C.}

    \UPSTIquestion{Aurait-on pu écrire ce programme en n'utilisant \textbf{que} des conditions enchaînées (\texttt{else if}), sans aucune imbrication, en combinant les tests avec \texttt{\&\&} ? Si oui, écrire la condition correspondant au tarif de \texttt{3} pour une voiture. Quelle version vous semble la plus lisible ?}
\end{UPSTIexercice}

\begin{UPSTIprofOnlyEnv}
    \begin{UPSTIcorrectionP}{Tarification d'un parking (exercice de synthèse)}
        Pseudo-code :
        \begin{verbatim}
SI type == 'M' ALORS
    SI dureeMinutes <= 15 ALORS
        AFFICHER "Tarif : 0"
    SINON
        AFFICHER "Tarif : 2"
    FIN SI
SINON SI type == 'V' ALORS
    SI dureeMinutes <= 15 ALORS
        AFFICHER "Tarif : 0"
    SINON SI dureeMinutes <= 60 ALORS
        AFFICHER "Tarif : 3"
    SINON
        AFFICHER "Tarif : 6"
    FIN SI
SINON
    AFFICHER "Type de vehicule inconnu"
FIN SI
        \end{verbatim}

        Traduction en C :
        \begin{lstlisting}[language=C]
if (type == 'M')
{
    if (dureeMinutes <= 15)
    {
        printf("Tarif : 0\n");
    }
    else
    {
        printf("Tarif : 2\n");
    }
}
else if (type == 'V')
{
    if (dureeMinutes <= 15)
    {
        printf("Tarif : 0\n");
    }
    else if (dureeMinutes <= 60)
    {
        printf("Tarif : 3\n");
    }
    else
    {
        printf("Tarif : 6\n");
    }
}
else
{
    printf("Type de vehicule inconnu\n");
}
        \end{lstlisting}

        Oui, c'est possible sans imbrication, à condition de répéter le test sur \texttt{type} dans chaque condition, exactement comme dans l'exercice \og Imbriqué ou enchaîné \fg. La condition du tarif de 3 pour une voiture deviendrait :
        \begin{lstlisting}[language=C]
else if (type == 'V' && dureeMinutes > 15 && dureeMinutes <= 60)
{
    printf("Tarif : 3\n");
}
        \end{lstlisting}
        La version imbriquée est ici plus lisible : elle regroupe naturellement toutes les règles propres à un même type de véhicule sous un seul test de \texttt{type}, évitant de le répéter dans chaque condition et réduisant le risque d'erreur si l'on doit modifier une règle plus tard.
    \end{UPSTIcorrectionP}
\end{UPSTIprofOnlyEnv}
